\section{Related Work}
\label{sec:related}

\subsection{AI-Assisted Creative Systems}

Research on AI-assisted creative generation has overwhelmingly focused on single-artifact
outputs: a poem, an image, a game level, a short narrative. Systems in this tradition treat
generation as a single forward pass from prompt to artifact, with quality evaluated against
criteria defined for that artifact class in isolation. Roemmele and Gordon
\cite{roemmele2015} demonstrate neural story continuation at the level of individual
narrative segments; Summerville et al.\ \cite{summerville2018} survey level generation
systems that produce one playable map at a time. The AI co-authorship literature
\cite{kreminski2020} similarly takes the individual artifact---a text, a scene, a game
element---as the unit of generation and the unit of evaluation.

Puzzle hunts, tabletop campaigns, and other multi-artifact creative systems present a
structurally different problem. The unit of generation is not a single artifact but a
\emph{coherent collection} of artifacts that must satisfy both per-artifact quality criteria
and cross-artifact consistency constraints: thematic unity, answer set compatibility,
difficulty arc, meta-puzzle closure. Generating such systems requires coordinated decisions
across multiple independent generation events, separated in time and context. Park et al.'s
generative agents \cite{park2023generative} demonstrate that LLMs can sustain coherent
behavior across multi-turn interactions through memory and identity representations, but
their setting is social simulation rather than structured creative production---the agents'
outputs are not artifacts consumed as inputs by downstream generation stages.

The key gap is that no prior work addresses coherent \emph{multi-artifact system generation}
across multiple independent generation runs, where each run's output becomes a structured
constraint on subsequent runs. Single-artifact systems cannot be trivially composed into
multi-artifact pipelines: coherence drift accumulates across generations unless explicitly
suppressed by a cross-stage consistency mechanism. Our pipeline addresses this gap through
stage-gated generation with a canonical world document that functions as a shared
read-only constraint across all stages.

\subsection{Procedural Content Generation with Constraints}

Procedural content generation (PCG) provides the nearest prior framework for understanding
structured creative pipelines. Shaker et al.\ \cite{shaker2016} survey the field broadly,
distinguishing generate-and-test approaches, constructive algorithms, and search-based
methods. Smith and Whitehead \cite{smith2010} develop Tanagra, a mixed-initiative level
design system that treats design constraints as real-time feedback to the generator,
demonstrating that constraint propagation across design decisions improves both output
quality and designer agency. The PCG-via-search paradigm in particular operationalizes
quality as a fitness function and treats generation as search over an artifact space
bounded by explicit constraints.

Our pipeline is constraint-based in the same spirit, but the constraint structure is
sequential across stages rather than concurrent within a single generation pass. Each
stage's output artifacts become binding constraints on subsequent stages: the world document
(Stage 1) constrains all thematic content generated in later stages; the puzzle pool ranking
(Stage 3) constrains which puzzle designs proceed to authoring; the feeder answers committed
at authoring time (Stage 6) constrain meta-puzzle design (Stage 5). This staged constraint
propagation differs from mixed-initiative PCG \cite{yannakakis2011} in an important
respect: the constraints are not relaxed or negotiated during generation, but locked at
explicit gate conditions. The world lock mechanism is not a soft preference but a hard
read-only boundary---downstream stages are explicitly prohibited from modifying it. This
design choice prioritizes coherence over flexibility, accepting that late-stage discoveries
cannot propagate back to world-level decisions.

The review-gate pattern that enforces these constraints---requiring an explicit pass verdict
from an evaluation stage before proceeding to the next generation stage---has no direct
precedent in the PCG literature, where quality filters are typically applied
post-hoc or embedded as fitness functions within a generative loop. Our gates are
pre-hoc: they prevent low-quality artifacts from entering the pipeline segment where
downstream authoring investment would be wasted on them.

\subsection{Puzzle Design and Quality}

The practitioner community has developed explicit quality standards for puzzle design that
are more codified than most creative domains. The Riven Standard, articulated by Rand Miller
in reflecting on the \emph{Myst} sequel, holds that a well-designed puzzle is one whose
mechanic is inseparable from its setting: the puzzle should \emph{be} what the domain does,
not merely occur within it \cite{miller2016riven}. Snyder \cite{snyder2019craftsmanship}
enumerates craftsmanship criteria for puzzle construction that translate across types: every
element should earn its presence, solutions should be uniquely determined, and the solver's
confirmation should be unambiguous. Competitive puzzle construction standards from the World
Puzzle Championship and U.S.\ Puzzle Championship instantiate these principles as
elimination criteria, providing a community-validated quality bar that is unusually explicit
relative to other creative fields.

Puzzle hunts introduce quality dimensions not present in standalone puzzle evaluation. A
puzzle hunt is not a collection of independent puzzles but a structured system: puzzles must
yield answer words that interact correctly with the meta-puzzle, maintain consistent
difficulty calibration across rounds, and sustain thematic coherence with the hunt's
fictional premise. The meta-puzzle itself introduces a structural correctness criterion---it
must be solvable from the feeder answers it receives, and it must close cleanly---that has
no analog in standalone puzzle design. Answer confirmation, a dimension in our Stage 7
rubric, is unique to the hunt context: it asks whether solving the puzzle produces an
unambiguous signal of correctness, enabling teams to verify their answer without
adjudication.

The AI expert panel review system deployed at Stages 3 and 7 of our pipeline operationalizes
these practitioner criteria as a six-dimension rubric (Clarity, Solvability, Elegance,
Reading Reward, Fun, Confirmation) with empirically grounded thresholds. The design and
calibration of this system is the subject of a companion paper
\cite{dellaLibera2025panel}; the present paper treats it as a component and characterizes
its behavior as a pipeline gate across six independent hunt scenarios.

\subsection{AI Pipelines for Complex Tasks}

The agent orchestration literature has proposed several frameworks for decomposing complex
tasks across multiple LLM invocations. Park et al.\ \cite{park2023generative} demonstrate
that generative agents with persistent memory and identity can sustain coherent behavior
across multi-turn social simulations. Nakajima's BabyAGI \cite{nakajima2023babyagi} and
Richards's AutoGPT \cite{richards2023autogpt} popularized autonomous agent loop
architectures in which a planning agent generates, prioritizes, and executes subtasks
without human intervention between steps. More structured orchestration frameworks such as
AutoGen \cite{autogen} and LangChain \cite{langchain} provide scaffolding for multi-agent
conversation and sequential tool use, enabling pipelines where agent outputs are routed as
inputs to downstream agents.

Our pipeline differs from this tradition in a design decision that is deliberate rather than
incidental: we do \emph{not} use autonomous agent loops. Each of the eleven stages is
human-triggerable and independently reviewable; no stage executes automatically as a
consequence of the previous stage completing. This design makes the pipeline slower than an
autonomous loop architecture, but it addresses a failure mode that autonomous loops inherit:
unchecked error propagation. In an autonomous loop, a generation error at Stage 3 that
escapes detection will cause downstream stages to invest authoring time on a flawed
artifact, and the error may not be visible until Stage 7 or Stage 10. In our pipeline, a
human reviews the output of each stage before triggering the next, and formal gate
conditions at Stages 3 and 7 require an explicit pass verdict before downstream work begins.

This \emph{review-gate pattern} relates to two bodies of practice outside the AI literature.
In traditional publishing, editorial pipelines impose stage gates between acquisitions,
developmental editing, copy editing, and production, with each gate requiring editorial
approval before investment in the next stage \cite{schiffrin2010}. In software engineering,
quality gates in continuous integration pipelines enforce structural analogs: automated and
human checks must pass before a change is promoted to the next environment. Our pipeline
adapts this logic to AI content generation, treating review stages not as overhead but as
the primary mechanism by which generation quality is enforced. The empirical question this
paper addresses is where in an eleven-stage pipeline these gates add the most value---and
the answer, across six independent hunt scenarios, is that Stages 3 and 7 account for the
majority of quality control, with failure rates at other stages negligible by comparison.
